\begin{center}
	\large{\textbf{\textcolor{black}{Desarrollo de un sistema embebido de adquisición de datos y telemetría satelital para boya instrumental EAMMRA.}}}\\
	\vspace{1pc}
	\textcolor {black}{Patricio BOS}
\end{center}

\bigskip
	\begin{center}
		\textcolor{black}{\textbf{RESUMEN}}	
	\end{center}
	\begin*
		\indent
	\end*
\textit{En este informe técnico se documenta el desarrollo de un sistema embebido de control, adquisición, supervisión, almacenamiento y telemetría satelital para la boya instrumental EAMMRA, en el marco del PIDDEF 0322. El sistema está orientado a la operación autónoma de una plataforma marítima capaz de adquirir variables ambientales, meteorológicas, de posición, energéticas y acústicas, almacenar los datos en forma local y transmitir información seleccionada mediante Iridium SBD. La implementación se basa en una computadora principal que ejecuta software desarrollado en Python bajo GNU/Linux. La arquitectura se organiza en módulos funcionales independientes, cada uno dividido en una capa de bajo nivel y una Máquina de Estados Finitos, coordinados por un proceso principal con enrutador de eventos y planificador de tareas central. Se describen la plataforma de hardware integrada, la arquitectura de software, las interfaces de comunicación, los formatos de datos, los mensajes binarios de telemetría y los mecanismos de robustez y diagnóstico. La verificación incluyó una \textit{suite} automatizada sin hardware, con 115 pruebas aprobadas y una cobertura global de líneas del 75~\%, además de ensayos de integración en laboratorio. El primer ensayo sostuvo la operación durante 118,7~h y permitió evaluar continuidad funcional, procesamiento acústico, solicitudes de telemetría y aislamiento frente a una falla inducida del bus USB--I2C. El segundo ensayo incorporó el subsistema panel solar--batería durante 101,8~h, registrando variaciones de tensión fotovoltaica, tensión de batería y estado de carga mediante el controlador XTRA2210. También se verificó la cadena de telemetría mediante solicitudes locales y recepción remota de mensajes Iridium SBD. Los resultados constituyen una validación funcional de integración en laboratorio y establecen una base técnica para mantenimiento, evolución del sistema y futuros ensayos de autonomía en condiciones de campo.}



\bigskip
	\begin{center}
		\textcolor{black}{\textbf{ABSTRACT}}
		\end{center}
			\begin*
			\indent
\textit{This technical report documents the development of an embedded system for control, data acquisition, monitoring, local storage, and satellite telemetry for the EAMMRA instrumented buoy, as part of PIDDEF 0322. The system is intended for autonomous operation on a maritime platform capable of acquiring environmental, meteorological, position, power-system, and acoustic data, storing them locally, and transmitting selected information through Iridium SBD. The implementation is based on a main onboard computer running Python software under GNU/Linux. The architecture is organized into independent functional modules, each divided into a low-level layer and a Finite State Machine, coordinated by a main process with an event router and a central task scheduler. The report describes the integrated hardware platform, software architecture, communication interfaces, data formats, binary telemetry messages, and robustness and diagnostic mechanisms. Verification included an automated no-hardware test suite, with 115 passed tests and 75~\% overall line coverage, as well as laboratory integration trials. The first trial sustained operation for 118.7~h and assessed functional continuity, acoustic processing, telemetry requests, and isolation after an induced USB--I2C bus failure. The second trial incorporated the solar panel--battery subsystem for 101.8~h, recording photovoltaic voltage, battery voltage, and state-of-charge variations through the XTRA2210 controller. The telemetry chain was also verified through local transmission requests and remote reception of Iridium SBD messages. The results provide a functional laboratory integration validation and establish a technical basis for maintenance, system evolution, and future autonomy trials under field conditions.}   
	\end*

\clearpage

